\documentclass[11pt,a4paper]{article}

% ── Packages ──
\usepackage[utf8]{inputenc}
\usepackage{amsmath,amssymb,amsthm,mathtools}
\usepackage{bm}          % bold math
\usepackage{geometry}
\usepackage{hyperref}
\usepackage{enumitem}
\usepackage{booktabs}
\usepackage{xcolor}
\usepackage{stmaryrd}     % for \llbracket, \rrbracket
\usepackage{tcolorbox}

\geometry{margin=1in}
\hypersetup{colorlinks=true, linkcolor=blue!60!black, citecolor=green!50!black}

% ── Theorem environments ──
\newtheorem{proposition}{Proposition}[section]
\newtheorem{lemma}[proposition]{Lemma}
\newtheorem{theorem}[proposition]{Theorem}
\theoremstyle{definition}
\newtheorem{definition}[proposition]{Definition}
\theoremstyle{remark}
\newtheorem{remark}[proposition]{Remark}

% ── Shortcuts ──
\newcommand{\R}{\mathbb{R}}
\newcommand{\dd}{\,\mathrm{d}}
\newcommand{\dt}{\partial_t}
\newcommand{\dn}{\partial_n}
\newcommand{\dv}{\operatorname{div}}
\newcommand{\curl}{\operatorname{curl}}
\newcommand{\grad}{\nabla}
\newcommand{\eps}{\varepsilon}
\newcommand{\half}{\tfrac{1}{2}}
\newcommand{\inner}[2]{\left\langle #1, #2 \right\rangle}
\newcommand{\norm}[1]{\left\| #1 \right\|}
\newcommand{\jump}[1]{\llbracket #1 \rrbracket}
\newcommand{\avg}[1]{\left\{ #1 \right\}}
\newcommand{\bfu}{\bm{u}}
\newcommand{\bfm}{\bm{m}}
\newcommand{\bfh}{\bm{h}}
\newcommand{\bfn}{\bm{n}}
\newcommand{\bff}{\bm{f}}
\newcommand{\bfz}{\bm{z}}
\newcommand{\bfv}{\bm{v}}
\newcommand{\bfM}{\bm{M}}
\newcommand{\bfH}{\bm{H}}
\newcommand{\bfB}{\bm{B}}
\newcommand{\Th}{\mathcal{T}_h}
\newcommand{\Eh}{\mathcal{E}_h}
\newcommand{\bfV}{\bm{V}}

% ── Colored remark boxes ──
\newtcolorbox{keypoint}{colback=blue!5!white, colframe=blue!50!black,
  fonttitle=\bfseries, title=Key Point}
\newtcolorbox{derivnote}{colback=yellow!5!white, colframe=orange!50!black,
  fonttitle=\bfseries, title=Derivation Note}

% ══════════════════════════════════════════════════════════════════════
\title{%
  Two-Phase Ferrohydrodynamics:\\
  Derivation from First Principles\\[6pt]
  \large Independent Formulation for Code Verification}
\author{Derivation Document — Semi\_Coupled Project}
\date{March 2026}

\begin{document}
\maketitle
\tableofcontents
\newpage

% ══════════════════════════════════════════════════════════════════════
% SECTION 1: CAHN-HILLIARD
% ══════════════════════════════════════════════════════════════════════
\section{Phase-Field Model: Cahn--Hilliard Equation}
\label{sec:CH}

\subsection{Physical Motivation}

We have two immiscible fluids in a bounded domain $\Omega \subset \R^d$
($d=2,3$): a ferrofluid and a non-magnetic carrier fluid. Rather than
tracking the sharp interface between them (which is numerically painful
and topologically fragile), we introduce a \emph{phase-field} variable
$\theta : \Omega \times [0,T] \to \R$ with the convention
\[
  \theta \approx +1 \quad \text{(ferrofluid)}, \qquad
  \theta \approx -1 \quad \text{(non-magnetic fluid)}.
\]
The interface is the thin transition layer where $\theta$ varies from
$-1$ to $+1$ over a width $\sim \eps$.

\subsection{Free Energy}

The Ginzburg--Landau free energy for the two-phase mixture is
\begin{equation}
  \label{eq:GL-energy}
  E_{\mathrm{GL}}[\theta]
  = \int_\Omega \frac{\eps}{2} |\grad\theta|^2
    + \frac{1}{\eps} F(\theta) \dd x,
\end{equation}
where $\eps > 0$ is the interface thickness parameter, and $F$ is a
double-well potential with minima at $\theta = \pm 1$.

\begin{derivnote}
Why this form? The term $\frac{\eps}{2}|\grad\theta|^2$ penalizes
spatial gradients (keeps the interface smooth), while
$\frac{1}{\eps}F(\theta)$ forces $\theta$ toward $\pm 1$ in the bulk.
The competition between these two terms produces a diffuse interface of
width $\mathcal{O}(\eps)$. As $\eps \to 0$, we formally recover a sharp
interface.
\end{derivnote}

\subsubsection{Choice of Double-Well Potential}

The standard polynomial double-well is $F_0(\theta) = \frac{1}{4}(\theta^2 - 1)^2$.
However, we need bounded derivatives for the analysis to work, so we use
a \emph{truncated} version:
\begin{equation}
  \label{eq:truncated-F}
  F(\theta) =
  \begin{cases}
    (\theta + 1)^2   & \text{if } \theta \leq -1, \\
    \frac{1}{4}(\theta^2 - 1)^2 & \text{if } -1 \leq \theta \leq 1, \\
    (\theta - 1)^2   & \text{if } \theta \geq 1.
  \end{cases}
\end{equation}

Set $f(\theta) := F'(\theta)$. Then:
\begin{itemize}
  \item In $[-1,1]$: $f(\theta) = \theta^3 - \theta$.
  \item For $\theta \leq -1$: $f(\theta) = 2(\theta+1)$.
  \item For $\theta \geq 1$: $f(\theta) = 2(\theta-1)$.
\end{itemize}

The crucial property we need is:
\begin{equation}
  \label{eq:f-bounds}
  |f(\theta)| \leq 2|\theta| + 1, \qquad
  |f'(\theta)| \leq 2, \qquad \forall\, \theta \in \R.
\end{equation}
The bound $|f'| \leq 2$ is what makes the convex-concave splitting work
later.

\subsection{Deriving the Cahn--Hilliard Equation}

The phase field should evolve to minimize the free energy while
conserving mass (total amount of each fluid is fixed). The standard
approach:

\paragraph{Step 1: Chemical potential.}
Define the chemical potential as the variational derivative of the free
energy:
\begin{equation}
  \label{eq:chem-potential}
  \psi := \frac{\delta E_{\mathrm{GL}}}{\delta \theta}
  = -\eps\,\Delta\theta + \frac{1}{\eps} f(\theta).
\end{equation}
To see this, take a variation $\theta \mapsto \theta + \delta\alpha\,\eta$
for test function $\eta$ with $\dn\theta = 0$ on $\partial\Omega$:
\begin{align*}
  \frac{d}{d\alpha}\bigg|_{\alpha=0} E_{\mathrm{GL}}[\theta + \alpha\eta]
  &= \int_\Omega \eps\,\grad\theta \cdot \grad\eta
     + \frac{1}{\eps} f(\theta)\,\eta \dd x \\
  &= \int_\Omega \left( -\eps\,\Delta\theta
     + \frac{1}{\eps} f(\theta) \right) \eta \dd x
     + \underbrace{\int_{\partial\Omega} \eps\,\dn\theta\,\eta \dd s}_{=\,0}.
\end{align*}

\paragraph{Step 2: Mass-conserving dynamics.}
To conserve mass ($\frac{d}{dt}\int_\Omega \theta \dd x = 0$), we need
$\dt\theta = -\dv \bm{J}$ for some flux $\bm{J}$. The simplest choice
that dissipates energy is $\bm{J} = -\gamma \grad\psi$ (Fick's law with
mobility $\gamma > 0$):
\begin{equation}
  \label{eq:CH}
  \dt\theta = \gamma\,\Delta\psi = -\gamma\,\Delta\!\left(
    -\eps\,\Delta\theta + \frac{1}{\eps} f(\theta)
  \right).
\end{equation}

\paragraph{Step 3: Verify energy dissipation.}
Multiply \eqref{eq:CH} by $\psi$ and integrate:
\begin{align}
  \frac{d}{dt} E_{\mathrm{GL}}[\theta]
  &= \int_\Omega \psi\,\dt\theta \dd x
  = \int_\Omega \psi\,\gamma\,\Delta\psi \dd x
  = -\gamma \int_\Omega |\grad\psi|^2 \dd x \leq 0.
  \label{eq:CH-dissipation}
\end{align}
Good: energy decreases.

\paragraph{Boundary conditions.}
Natural (Neumann) on both $\theta$ and $\psi$:
\begin{equation}
  \label{eq:CH-BC}
  \dn\theta = 0, \qquad \dn\psi = 0 \qquad \text{on } \partial\Omega.
\end{equation}

\subsection{First-Order System Form}

For numerics, we write CH as a system of two second-order PDEs:
\begin{equation}
  \label{eq:CH-system}
  \boxed{
  \begin{aligned}
    \dt\theta + \bfu \cdot \grad\theta &= \gamma\,\Delta\psi, \\
    \psi &= -\eps\,\Delta\theta + \frac{1}{\eps} f(\theta),
  \end{aligned}
  }
\end{equation}
where we have added the advection term $\bfu \cdot \grad\theta$ because
the phase field is carried by the fluid velocity $\bfu$ (this will be
justified in Section~\ref{sec:NS}).

\begin{remark}
  Mass conservation: integrating the first equation over $\Omega$ with
  $\dn\psi = 0$ and $\dv\bfu = 0$ gives
  $\frac{d}{dt}\int_\Omega \theta \dd x = 0$.
  The advection term contributes zero because
  $\int_\Omega \bfu \cdot \grad\theta \dd x
   = -\int_\Omega (\dv\bfu)\,\theta \dd x = 0$
  for divergence-free $\bfu$.
\end{remark}


% ══════════════════════════════════════════════════════════════════════
% SECTION 2: FERROHYDRODYNAMICS — SHLIOMIS
% ══════════════════════════════════════════════════════════════════════
\section{Magnetization Dynamics: Shliomis Model}
\label{sec:shliomis}

\subsection{Physical Setting}

A ferrofluid is a colloidal suspension of magnetic nanoparticles
($\sim 10$\,nm diameter) in a carrier liquid. Each particle has a
permanent magnetic dipole moment. The \emph{magnetization}
$\bfm : \Omega \times [0,T] \to \R^d$ is the volume-averaged magnetic
moment density.

When an external magnetic field $\bfh$ is applied, the particles tend to
align with $\bfh$, but thermal fluctuations and viscous drag resist
alignment. The \emph{Shliomis model} describes this competition.

\subsection{Full Shliomis Model (Single Phase)}

The full monophase Shliomis model couples Navier--Stokes with a
magnetization evolution equation. The magnetization equation is:
\begin{equation}
  \label{eq:shliomis-full}
  \dt\bfm + (\bfu \cdot \grad)\bfm - \half\,\curl(\bfu) \times \bfm
  = -\frac{1}{\tau}\left(\bfm - \chi_0\,\bfh\right)
    - \beta\,\bfm \times (\bfm \times \bfh),
\end{equation}
where:
\begin{itemize}[nosep]
  \item $\bfu$: fluid velocity,
  \item $\tau > 0$: magnetic relaxation time,
  \item $\chi_0 > 0$: magnetic susceptibility (equilibrium: $\bfm = \chi_0\,\bfh$),
  \item $\beta \geq 0$: coefficient for the cross-relaxation term,
  \item $\half\,\curl(\bfu) \times \bfm$: vorticity coupling (particles
        rotate with the fluid).
\end{itemize}

\subsection{Simplification: Why We Drop Terms}

For commercial ferrofluids, $\tau \sim 10^{-5}$ to $10^{-9}$\,s, so
$1/\tau$ is enormous. This means:

\begin{enumerate}
  \item The relaxation term $-\frac{1}{\tau}(\bfm - \chi_0\bfh)$ dominates.
        On time scales of interest (flow time $\sim \mathcal{O}(1)$\,s),
        the magnetization is always close to equilibrium:
        $\bfm \approx \chi_0\bfh$.

  \item The cross-relaxation term $\beta\,\bfm \times (\bfm \times \bfh)$
        is a second-order correction. Near equilibrium,
        $\bfm \times \bfh \approx 0$, so this term is negligible.

  \item The vorticity coupling $\half\curl(\bfu) \times \bfm$:
        this term cannot be handled in a clean energy framework because
        it does not have a symmetric structure. Dropping it simplifies
        the energy analysis.
\end{enumerate}

\begin{keypoint}
The simplified Shliomis magnetization equation is:
\begin{equation}
  \label{eq:shliomis-simple}
  \boxed{
    \dt\bfm + (\bfu \cdot \grad)\bfm + \frac{1}{\tau}\bfm
    = \frac{\chi_0}{\tau}\bfh.
  }
\end{equation}
This is a \emph{linear advection-reaction equation} for $\bfm$ given
$\bfu$ and $\bfh$. The $(1/\tau)\bfm$ term provides damping, and
$(\chi_0/\tau)\bfh$ drives $\bfm$ toward $\chi_0\bfh$.
\end{keypoint}

\begin{derivnote}
Even though $\tau$ is small and $\bfm \approx \chi_0\bfh$ in steady
state, we do \emph{not} replace the PDE by the algebraic relation
$\bfm = \chi_0\bfh$. The PDE form is essential because:
(a) it provides the correct energy structure for stability,
(b) it allows for transient dynamics during field changes,
(c) it naturally handles the two-phase case where $\chi_0$ varies in space.
\end{derivnote}

\subsection{Energy Structure of the Magnetization Equation}

Multiply \eqref{eq:shliomis-simple} by $\bfm$ and integrate over $\Omega$.
For the advection term, assuming $\dv\bfu = 0$ and $\bfu \cdot \bfn = 0$
on $\partial\Omega$:
\[
  \int_\Omega [(\bfu \cdot \grad)\bfm] \cdot \bfm \dd x
  = \half \int_\Omega \bfu \cdot \grad|\bfm|^2 \dd x
  = -\half \int_\Omega (\dv\bfu)|\bfm|^2 \dd x = 0.
\]
So:
\begin{equation}
  \label{eq:mag-energy-test}
  \half\frac{d}{dt}\norm{\bfm}^2_{L^2}
  + \frac{1}{\tau}\norm{\bfm}^2_{L^2}
  = \frac{\chi_0}{\tau}\int_\Omega \bfh \cdot \bfm \dd x.
\end{equation}
The right-hand side couples to the magnetic field energy — we will close
this in Section~\ref{sec:energy}.


% ══════════════════════════════════════════════════════════════════════
% SECTION 3: MAGNETOSTATICS
% ══════════════════════════════════════════════════════════════════════
\section{Magnetostatics: From Maxwell to Poisson}
\label{sec:maxwell}

\subsection{Maxwell's Equations in Matter (Magnetostatics)}

We assume the magnetic field varies slowly compared to the speed of
light (no electromagnetic waves), so we use the \emph{magnetostatic}
approximation. The relevant Maxwell equations are:
\begin{equation}
  \label{eq:maxwell}
  \curl(\bfh) = \bm{0}, \qquad
  \dv(\bfB) = 0,
\end{equation}
where $\bfB$ is the magnetic flux density and $\bfh$ is the magnetic
field intensity. The constitutive relation in a magnetic medium is:
\begin{equation}
  \label{eq:constitutive}
  \bfB = \mu_0(\bfh + \bfm),
\end{equation}
with $\mu_0$ the permeability of free space and $\bfm$ the magnetization.

\subsection{Deriving the Poisson Equation}

\paragraph{Step 1:} Since $\curl(\bfh) = \bm{0}$ in a simply connected
domain, there exists a scalar potential $\varphi$ such that
\begin{equation}
  \label{eq:h-grad-phi}
  \bfh = -\grad\varphi + \bfh_a,
\end{equation}
where $\bfh_a$ is the externally applied field (given data). We decompose:
\[
  \bfh = \bfh_a + \bfh_d, \qquad \bfh_d = -\grad\varphi,
\]
where $\bfh_d$ is the \emph{demagnetizing field} (the field created by
the magnetic material itself).

\begin{derivnote}
Why this decomposition? The applied field $\bfh_a$ satisfies
$\curl(\bfh_a) = \bm{0}$ and $\dv(\bfh_a) = 0$ in free space (it is
the field that would exist without the ferrofluid). The demagnetizing
field $\bfh_d$ is generated by the magnetized material. In our
formulation, $\bfh_a$ is given and we solve for $\varphi$ (which gives
$\bfh_d$).
\end{derivnote}

\paragraph{Step 2:} Substitute into $\dv(\bfB) = 0$:
\begin{align*}
  0 &= \dv(\bfB) = \mu_0\,\dv(\bfh + \bfm) \\
    &= \mu_0\,\dv(\bfh_a + \bfh_d + \bfm) \\
    &= \mu_0\left(\underbrace{\dv(\bfh_a)}_{=\,0}
       + \dv(-\grad\varphi) + \dv(\bfm)\right) \\
    &= \mu_0\left(-\Delta\varphi + \dv(\bfm)\right).
\end{align*}

Therefore:
\begin{equation}
  \label{eq:poisson}
  \boxed{
    -\Delta\varphi = -\dv(\bfm) \qquad \text{in } \Omega.
  }
\end{equation}
Or equivalently: $\Delta\varphi = \dv(\bfm)$.

\subsection{Boundary Conditions}

On $\partial\Omega$, we impose:
\begin{equation}
  \label{eq:poisson-BC}
  \dn\varphi = \bfm \cdot \bfn \qquad \text{on } \partial\Omega.
\end{equation}

\begin{derivnote}
This comes from the physical requirement that
$\bfB \cdot \bfn$ is continuous across $\partial\Omega$. Outside
$\Omega$, $\bfm = \bm{0}$, so $\bfB = \mu_0\bfh$. Continuity of
$\bfB \cdot \bfn$ gives
$(\bfh_a + \bfh_d + \bfm) \cdot \bfn |_{\text{inside}}
= (\bfh_a + \bfh_d) \cdot \bfn|_{\text{outside}}$,
which simplifies to $\bfh_d \cdot \bfn = -\bfm \cdot \bfn$, i.e.,
$-\dn\varphi = -\bfm \cdot \bfn$, giving \eqref{eq:poisson-BC}.
\end{derivnote}

\subsection{Total Magnetic Field}

After solving the Poisson equation, the total field is:
\begin{equation}
  \label{eq:total-h}
  \bfh = \bfh_a - \grad\varphi.
\end{equation}

\begin{remark}
  The magnetization $\bfm$ and field $\bfh$ are coupled:
  $\bfm$ depends on $\bfh$ through \eqref{eq:shliomis-simple}, and
  $\bfh$ depends on $\bfm$ through \eqref{eq:poisson}. This nonlinear
  coupling is the reason we need Picard iteration later.
\end{remark}


% ══════════════════════════════════════════════════════════════════════
% SECTION 4: NAVIER-STOKES WITH BODY FORCES
% ══════════════════════════════════════════════════════════════════════
\section{Navier--Stokes with Magnetic and Capillary Forces}
\label{sec:NS}

\subsection{Incompressible Navier--Stokes}

The fluid velocity $\bfu$ and pressure $p$ satisfy the incompressible
Navier--Stokes equations:
\begin{equation}
  \label{eq:NS-generic}
  \begin{aligned}
    \rho\left(\dt\bfu + (\bfu \cdot \grad)\bfu\right)
    - \dv\!\left(2\nu\,\bm{D}(\bfu)\right) + \grad p &= \bff, \\
    \dv(\bfu) &= 0,
  \end{aligned}
\end{equation}
where $\bm{D}(\bfu) = \half(\grad\bfu + \grad\bfu^\top)$ is the
symmetric strain rate, $\nu$ is the dynamic viscosity, $\rho$ is the
density, and $\bff$ collects all body forces.

\subsection{Kelvin Force (Magnetic Body Force)}
\label{sec:kelvin}

The force exerted by a magnetic field on a magnetized material is the
\emph{Kelvin force}:
\begin{equation}
  \label{eq:kelvin}
  \bff_{\mathrm{mag}} = \mu_0\,(\bfm \cdot \grad)\bfh.
\end{equation}

\begin{derivnote}
Derivation: Consider a magnetic dipole $\bfm$ in a non-uniform field
$\bfh$. Its potential energy is $U = -\mu_0\,\bfm \cdot \bfh$. The
force on the dipole is $\bff = -\grad U = \mu_0\,\grad(\bfm \cdot \bfh)$.
Using the vector identity
\[
  \grad(\bfm \cdot \bfh)
  = (\bfm \cdot \grad)\bfh + (\bfh \cdot \grad)\bfm
    + \bfm \times \curl(\bfh) + \bfh \times \curl(\bfm),
\]
and noting that $\curl(\bfh) = \bm{0}$ (magnetostatics), this becomes
\[
  \bff = \mu_0\left[(\bfm \cdot \grad)\bfh + (\bfh \cdot \grad)\bfm
         + \bfh \times \curl(\bfm)\right].
\]
For a continuum of dipoles, after averaging, the standard result is the
Kelvin force $\bff_{\mathrm{mag}} = \mu_0(\bfm \cdot \grad)\bfh$.
The other terms either cancel or are absorbed into a modified pressure.
\end{derivnote}

\begin{remark}
  In the full Shliomis model, there is also an antisymmetric stress
  contribution $\frac{\mu_0}{2}\curl(\bfm \times \bfh)$. This was
  dropped in the simplified model (Section~\ref{sec:shliomis}).
\end{remark}

\subsection{Capillary Force (Surface Tension)}
\label{sec:capillary}

In the diffuse-interface framework, surface tension appears as a
distributed body force derived from the capillary stress tensor.

The capillary stress tensor is
$\bm{\sigma}_c = \lambda\,\grad\theta \otimes \grad\theta$,
where $\lambda > 0$ is proportional to surface tension.
The capillary force is:
\begin{equation}
  \label{eq:capillary-force-div}
  \bff_c = \dv(\bm{\sigma}_c) = \lambda\left[
    \Delta\theta\,\grad\theta
    + \half\grad|\grad\theta|^2
  \right].
\end{equation}

Using the chemical potential $\psi$ from \eqref{eq:chem-potential}, we
can rewrite this more conveniently. Note that
\[
  \psi = -\eps\,\Delta\theta + \frac{1}{\eps}f(\theta)
  \quad \Longrightarrow \quad
  \Delta\theta = -\frac{1}{\eps}\psi + \frac{1}{\eps^2}f(\theta).
\]
After some manipulation (using $\grad F(\theta) = f(\theta)\grad\theta$),
the capillary force can be written as:
\begin{equation}
  \label{eq:capillary-force}
  \boxed{
    \bff_c = \frac{\lambda}{\eps}\,\theta\,\grad\psi.
  }
\end{equation}

\begin{derivnote}
The key trick: write
$\dv(\bm{\sigma}_c) = \lambda\Delta\theta\,\grad\theta + \text{(gradient terms)}$.
The gradient terms can be absorbed into a modified pressure
$\tilde{p} = p + \frac{\lambda}{2}|\grad\theta|^2 + \frac{\lambda}{\eps}F(\theta)$
(see the original derivation in the Cahn--Hilliard literature, e.g.,
Feng 2006, Kay et al.\ 2009). What remains is
$\lambda\Delta\theta\,\grad\theta$, which using
$\psi = -\eps\Delta\theta + \frac{1}{\eps}f(\theta)$ gives us the
$\frac{\lambda}{\eps}\theta\grad\psi$ form after absorbing remaining
gradient terms into $\tilde{p}$.
\end{derivnote}

\subsection{Gravity (Boussinesq Approximation)}

For the Rosensweig instability, gravity is essential (it opposes the
magnetic force that creates spikes). Using Boussinesq:
\begin{equation}
  \label{eq:gravity}
  \bff_g = -\rho_0\left(1 + \alpha_\theta\,\theta\right)\bm{g},
\end{equation}
where $\rho_0$ is a reference density, $\alpha_\theta$ accounts for the
density contrast, and $\bm{g}$ is gravitational acceleration.

\subsection{Complete Navier--Stokes}

Collecting all forces:
\begin{equation}
  \label{eq:NS-complete}
  \boxed{
  \begin{aligned}
    \dt\bfu + (\bfu \cdot \grad)\bfu
    - \dv\!\left(2\nu(\theta)\,\bm{D}(\bfu)\right)
    + \grad p
    &= \frac{\lambda}{\eps}\,\theta\,\grad\psi
       + \mu_0\,(\bfm \cdot \grad)\bfh
       + \bff_g, \\
    \dv(\bfu) &= 0,
  \end{aligned}
  }
\end{equation}
with boundary conditions $\bfu = \bm{0}$ on $\partial\Omega$ (no-slip).


% ══════════════════════════════════════════════════════════════════════
% SECTION 5: TWO-PHASE EXTENSION
% ══════════════════════════════════════════════════════════════════════
\section{Two-Phase Material Properties}
\label{sec:two-phase}

In the two-phase system, material properties depend on the phase field
$\theta$. Let subscript 1 denote ferrofluid ($\theta = +1$) and
subscript 2 denote non-magnetic fluid ($\theta = -1$).

\subsection{Interpolation Functions}

We use linear interpolation:
\begin{equation}
  \label{eq:interp}
  \begin{aligned}
    \nu(\theta) &= \frac{\nu_1 + \nu_2}{2}
                 + \frac{\nu_1 - \nu_2}{2}\,\theta, \\
    \chi_0(\theta) &= \frac{\chi_{0,1} + \chi_{0,2}}{2}
                    + \frac{\chi_{0,1} - \chi_{0,2}}{2}\,\theta.
  \end{aligned}
\end{equation}
Since the non-magnetic fluid has $\chi_{0,2} = 0$:
\begin{equation}
  \label{eq:chi-interp}
  \chi_0(\theta) = \frac{\chi_{0,1}}{2}(1 + \theta).
\end{equation}

\begin{remark}
  This means $\chi_0(\theta) > 0$ only where $\theta > -1$ (i.e., where
  ferrofluid is present). In the non-magnetic region ($\theta = -1$),
  $\chi_0 = 0$ and the magnetization relaxes to zero.
\end{remark}

\subsection{Relaxation Time}

Similarly for the relaxation time, but since the non-magnetic fluid has
no magnetic relaxation, we set $\tau(\theta) = \tau_1$ in the ferrofluid
region and handle the non-magnetic region through $\chi_0 = 0$
(the magnetization equation naturally drives $\bfm \to 0$ when
$\chi_0 = 0$).


% ══════════════════════════════════════════════════════════════════════
% SECTION 6: COMPLETE SYSTEM
% ══════════════════════════════════════════════════════════════════════
\section{Complete Two-Phase Ferrohydrodynamic System}
\label{sec:complete}

Collecting everything from the previous sections, the complete system in
$\Omega \times (0,T]$ is:

\bigskip
\begin{equation}
  \label{eq:system}
  \boxed{
  \begin{aligned}
    &\textbf{Cahn--Hilliard:} \\
    &\quad \dt\theta + \bfu \cdot \grad\theta
      = \gamma\,\Delta\psi, \\
    &\quad \psi = -\eps\,\Delta\theta + \frac{1}{\eps}f(\theta), \\[6pt]
    &\textbf{Magnetization:} \\
    &\quad \dt\bfm + (\bfu \cdot \grad)\bfm + \frac{1}{\tau}\bfm
      = \frac{\chi_0(\theta)}{\tau}\bfh, \\[6pt]
    &\textbf{Magnetostatics:} \\
    &\quad -\Delta\varphi = -\dv(\bfm), \qquad
      \bfh = \bfh_a - \grad\varphi, \\[6pt]
    &\textbf{Navier--Stokes:} \\
    &\quad \dt\bfu + (\bfu \cdot \grad)\bfu
      - \dv\!\left(2\nu(\theta)\,\bm{D}(\bfu)\right) + \grad p \\
    &\qquad\qquad = \frac{\lambda}{\eps}\,\theta\,\grad\psi
       + \mu_0\,(\bfm \cdot \grad)\bfh + \bff_g, \\
    &\quad \dv(\bfu) = 0.
  \end{aligned}
  }
\end{equation}

\bigskip\noindent
\textbf{Boundary conditions:}
\begin{equation}
  \label{eq:BCs}
  \dn\theta = 0, \quad \dn\psi = 0, \quad
  \bfu = \bm{0}, \quad
  \dn\varphi = \bfm \cdot \bfn
  \qquad \text{on } \partial\Omega.
\end{equation}


% ══════════════════════════════════════════════════════════════════════
% SECTION 7: ENERGY LAW
% ══════════════════════════════════════════════════════════════════════
\section{Energy Law for the Continuous System}
\label{sec:energy}

\subsection{Total Energy Functional}

Define the total energy:
\begin{equation}
  \label{eq:total-energy}
  E(t) = \underbrace{\frac{1}{2}\norm{\bfu}^2_{L^2}}_{\text{kinetic}}
       + \underbrace{\frac{\lambda}{\eps}\int_\Omega
         \frac{\eps}{2}|\grad\theta|^2
         + \frac{1}{\eps}F(\theta) \dd x}_{\text{Ginzburg--Landau}}
       + \underbrace{\frac{\mu_0}{2\chi_0}\norm{\bfm}^2_{L^2}
         + \frac{\mu_0}{2}\norm{\grad\varphi}^2_{L^2}}_{\text{magnetic}}.
\end{equation}

\begin{derivnote}
Why these specific coefficients?
\begin{itemize}[nosep]
  \item The kinetic energy $\frac{1}{2}\|\bfu\|^2$ is standard.
  \item The GL energy is multiplied by $\lambda/\eps$ because the
        capillary force involves this coefficient.
  \item The magnetic energy $\frac{\mu_0}{2\chi_0}\|\bfm\|^2$ comes
        from the magnetization equation: testing with
        $\frac{\mu_0}{\chi_0}\bfm$ produces
        $\frac{\mu_0}{2\chi_0}\frac{d}{dt}\|\bfm\|^2$.
  \item The field energy $\frac{\mu_0}{2}\|\grad\varphi\|^2$ comes
        from the Poisson equation: testing with $\mu_0\,\dt\varphi$.
\end{itemize}
The art is choosing test functions so that all coupling terms cancel.
\end{derivnote}

\subsection{Energy Identity}

\begin{proposition}[Energy Dissipation]
  \label{prop:energy}
  Smooth solutions of system~\eqref{eq:system} satisfy:
  \begin{equation}
    \label{eq:energy-law}
    \frac{d}{dt} E(t)
    + 2\int_\Omega \nu(\theta)\,|\bm{D}(\bfu)|^2 \dd x
    + \frac{\lambda\gamma}{\eps}\norm{\grad\psi}^2_{L^2}
    + \frac{\mu_0}{\tau\chi_0}\norm{\bfm}^2_{L^2}
    = \frac{\mu_0}{\tau}\int_\Omega \bfh_a \cdot \bfm \dd x.
  \end{equation}
\end{proposition}

\begin{proof}[Proof sketch]
Test each equation with the appropriate function and add:

\paragraph{Step 1: NS tested with $\bfu$.}
\[
  \half\frac{d}{dt}\norm{\bfu}^2
  + 2\int_\Omega \nu|\bm{D}(\bfu)|^2 \dd x
  = \frac{\lambda}{\eps}\int_\Omega \theta\,\grad\psi \cdot \bfu \dd x
    + \mu_0 \int_\Omega [(\bfm \cdot \grad)\bfh] \cdot \bfu \dd x.
\]
The pressure term vanishes: $\int \grad p \cdot \bfu = -\int p\,\dv\bfu = 0$.
The advection nonlinearity vanishes:
$\int [(\bfu\cdot\grad)\bfu]\cdot\bfu = 0$ (skew-symmetry for $\dv\bfu=0$).

\paragraph{Step 2: CH tested with $\frac{\lambda}{\eps}\psi$.}
Using $\dt\theta = -\bfu\cdot\grad\theta + \gamma\Delta\psi$:
\[
  \frac{\lambda}{\eps}\int_\Omega \psi\,\dt\theta \dd x
  = \frac{\lambda}{\eps}\frac{d}{dt}\left[
    \int_\Omega \frac{\eps}{2}|\grad\theta|^2 + \frac{1}{\eps}F(\theta) \dd x
  \right],
\]
which equals $\frac{d}{dt}E_{\mathrm{GL}}$. The cross-term with $\bfu$:
\[
  -\frac{\lambda}{\eps}\int_\Omega \psi\,\bfu\cdot\grad\theta \dd x.
\]
And the dissipation: $-\frac{\lambda\gamma}{\eps}\|\grad\psi\|^2$.

\paragraph{Step 3: Magnetization tested with $\frac{\mu_0}{\chi_0}\bfm$.}
\[
  \frac{\mu_0}{2\chi_0}\frac{d}{dt}\norm{\bfm}^2
  + \underbrace{\frac{\mu_0}{\chi_0}\int_\Omega [(\bfu\cdot\grad)\bfm]\cdot\bfm}_{=\,0}
  + \frac{\mu_0}{\tau\chi_0}\norm{\bfm}^2
  = \frac{\mu_0}{\tau}\int_\Omega \bfh \cdot \bfm \dd x.
\]

\paragraph{Step 4: Poisson tested with $\mu_0\,\dt\varphi$.}
From $-\Delta\varphi = -\dv\bfm$:
\[
  \mu_0\int_\Omega \grad\varphi \cdot \grad(\dt\varphi) \dd x
  = \mu_0\int_\Omega \bfm \cdot \grad(\dt\varphi) \dd x.
\]
The left side gives $\frac{\mu_0}{2}\frac{d}{dt}\|\grad\varphi\|^2$.

\paragraph{Step 5: Cancel cross-terms.}
The capillary coupling: $\frac{\lambda}{\eps}\int\theta\,\grad\psi\cdot\bfu$
from NS cancels with $-\frac{\lambda}{\eps}\int\psi\,\bfu\cdot\grad\theta$ from CH
because (integration by parts with $\dv\bfu = 0$):
\[
  \int_\Omega \theta\,\grad\psi\cdot\bfu \dd x
  = -\int_\Omega \psi\,\dv(\theta\bfu) \dd x
  = -\int_\Omega \psi\left(\bfu\cdot\grad\theta + \theta\,\dv\bfu\right) \dd x
  = -\int_\Omega \psi\,\bfu\cdot\grad\theta \dd x. \quad \checkmark
\]

The magnetic coupling: $\mu_0\int[(\bfm\cdot\grad)\bfh]\cdot\bfu$ from NS
and $\frac{\mu_0}{\tau}\int\bfh\cdot\bfm$ from magnetization
combine with the Poisson contribution. The key identity is:
\[
  \int_\Omega [(\bfm\cdot\grad)\bfh]\cdot\bfu \dd x
  = -\int_\Omega [(\bfu\cdot\grad)\bfm]\cdot\bfh \dd x
    - \int_\Omega [(\bfu\cdot\grad)\bfh]\cdot\bfm \dd x.
\]
Together with the Poisson time-derivative term, everything closes to
give \eqref{eq:energy-law}.
\end{proof}

\begin{keypoint}
The right-hand side $\frac{\mu_0}{\tau}\int\bfh_a\cdot\bfm$ is the
power input from the external field. In steady state, this power is
balanced by viscous dissipation, phase-field diffusion, and magnetic
relaxation.
\end{keypoint}

\begin{remark}
  \label{rem:chi-restriction}
  For the energy estimate to close properly in the discrete setting, one
  needs the demagnetizing field energy to be controlled. This leads to
  the restriction $\chi_0 \leq 4$ in practice (the constant comes from
  the Poincar\'e-type estimate for the Poisson problem). This is a
  limitation of the scheme, not the physics.
\end{remark}


% ══════════════════════════════════════════════════════════════════════
% SECTION 8: FINITE ELEMENT DISCRETIZATION
% ══════════════════════════════════════════════════════════════════════
\section{Finite Element Discretization}
\label{sec:fem}

\subsection{Mesh and Notation}

Let $\Th$ be a conforming quadrilateral (2D) or hexahedral (3D)
partition of $\Omega$, with mesh size $h$. Let $\Eh$ denote the set of
interior faces (edges in 2D). For each face $e \in \Eh$ shared by
elements $K^+$ and $K^-$:
\begin{itemize}[nosep]
  \item $\bfn_e$: unit normal (with fixed orientation),
  \item $\jump{v} := v^+ - v^-$: jump across $e$,
  \item $\avg{v} := \half(v^+ + v^-)$: average across $e$.
\end{itemize}

\subsection{Finite Element Spaces}
\label{sec:fe-spaces}

\paragraph{Why mixed CG/DG?}
The magnetization equation has dominant advection
($(\bfu\cdot\grad)\bfm$), so standard CG would produce oscillations.
DG handles advection naturally via upwinding. However, the Poisson
equation for $\varphi$ needs CG for efficiency (DG Poisson is expensive
due to the penalty/LDG terms).

The critical constraint linking these spaces is:
\begin{equation}
  \label{eq:space-constraint}
  \boxed{
    \grad(X_h) \subseteq \bfM_h,
  }
\end{equation}
where $X_h$ is the CG space for $\varphi$ and $\bfM_h$ is the DG space
for $\bfm$. This ensures that testing the Poisson equation is consistent
with the magnetization discretization.

\begin{derivnote}
Why is this constraint needed? In the energy proof, we test the Poisson
equation with $\dt\varphi_h$, which gives terms involving
$\grad(\dt\varphi_h)$. This gradient must live in $\bfM_h$ for the
discrete integration by parts to work correctly. If $\grad(X_h) \not\subseteq \bfM_h$, the discrete energy law breaks.
\end{derivnote}

\paragraph{Concrete spaces:}
\begin{equation}
  \label{eq:spaces}
  \begin{aligned}
    \text{Velocity:}\quad &\bfV{}_h \subset [H^1_0(\Omega)]^d,
      &&\text{CG } \bm{Q}_2, \\
    \text{Pressure:}\quad &Q_h \subset L^2(\Omega),
      &&\text{DG } P_1, \\
    \text{Phase field:}\quad &S_h \subset H^1(\Omega),
      &&\text{CG } Q_2, \\
    \text{Poisson:}\quad &X_h \subset H^1(\Omega),
      &&\text{CG } Q_2, \\
    \text{Magnetization:}\quad &\bfM_h \subset [L^2(\Omega)]^d,
      &&\text{DG } \bm{Q}_1.
  \end{aligned}
\end{equation}

\begin{keypoint}
$\varphi \in Q_2$ (CG) and $\bfm \in Q_1$ (DG). Since
$\grad(Q_2) \subset Q_1$ (the gradient of a quadratic polynomial is
linear), the constraint \eqref{eq:space-constraint} is satisfied.
If we used $Q_1$ for $\varphi$, we would need $Q_0$ for $\bfm$ (constants),
which is too low-order.
\end{keypoint}

\subsection{DG Transport Form for Magnetization}
\label{sec:dg-transport}

The advection term $(\bfu\cdot\grad)\bfm$ needs special treatment in
DG. We derive the DG bilinear form from scratch.

\paragraph{Step 1: Start from the weak form on each element.}
On element $K$, multiply $(\bfu\cdot\grad)\bfm$ by test function $\bfz$
and integrate:
\[
  \int_K [(\bfu\cdot\grad)\bfm] \cdot \bfz \dd x.
\]

\paragraph{Step 2: Integration by parts on $K$.}
\begin{align*}
  \int_K (\bfu\cdot\grad)\bfm \cdot \bfz \dd x
  &= -\int_K \bfm \cdot \dv(\bfu \otimes \bfz) \dd x
     + \int_{\partial K} (\bfu \cdot \bfn_K)\,\bfm \cdot \bfz \dd s \\
  &= -\int_K \bfm \cdot [(\bfu\cdot\grad)\bfz] \dd x
     - \int_K (\dv\bfu)\,\bfm \cdot \bfz \dd x
     + \int_{\partial K} (\bfu\cdot\bfn_K)\,\bfm \cdot \bfz \dd s.
\end{align*}

\paragraph{Step 3: Sum over all elements.}
\[
  \sum_K \int_K (\bfu\cdot\grad)\bfm\cdot\bfz \dd x
  = -\sum_K \int_K \bfm \cdot [(\bfu\cdot\grad)\bfz] \dd x
    - \sum_K \int_K (\dv\bfu)\bfm\cdot\bfz \dd x
    + \sum_K \int_{\partial K} (\bfu\cdot\bfn_K)\bfm\cdot\bfz \dd s.
\]
The boundary integral sum, after pairing faces, gives jump terms on
interior faces:
\[
  \sum_K \int_{\partial K} (\bfu\cdot\bfn_K)\bfm\cdot\bfz \dd s
  = \sum_{e \in \Eh} \int_e \left[
    \avg{\bfu\cdot\bfn}\,\jump{\bfm\cdot\bfz}
    + \jump{\bfu\cdot\bfn}\,\avg{\bfm\cdot\bfz}
  \right] \dd s.
\]

\paragraph{Step 4: Simplify using $\dv\bfu = 0$.}
For continuous velocity ($\bfu \in \bfV{}_h$, CG), $\jump{\bfu\cdot\bfn} = 0$.
Also $\dv\bfu = 0$ (pointwise for the discrete divergence-free velocity).
So:
\[
  b_h(\bfu, \bfm, \bfz) := \sum_K \int_K [(\bfu\cdot\grad)\bfm]\cdot\bfz \dd x
  = -\sum_K \int_K \bfm\cdot[(\bfu\cdot\grad)\bfz] \dd x
    + \sum_{e \in \Eh} \int_e (\bfu\cdot\bfn)\,\jump{\bfm\cdot\bfz} \dd s.
\]

\paragraph{Step 5: Symmetrize (skew-symmetric form).}
Define the DG trilinear form:
\begin{equation}
  \label{eq:DG-trilinear}
  \boxed{
  \begin{aligned}
    B_h(\bfu, \bfm, \bfz) &:=
    \sum_K \int_K [(\bfu\cdot\grad)\bfm]\cdot\bfz \dd x \\
    &\quad + \half \sum_K \int_K (\dv\bfu)\,\bfm\cdot\bfz \dd x \\
    &\quad - \half \sum_{e \in \Eh} \int_e
      (\bfu\cdot\bfn)\,\jump{\bfm}\cdot\avg{\bfz} \dd s \\
    &\quad - \half \sum_{e \in \Eh} \int_e
      (\bfu\cdot\bfn)\,\avg{\bfm}\cdot\jump{\bfz} \dd s.
  \end{aligned}
  }
\end{equation}

\begin{derivnote}
The skew-symmetrization is done so that
$B_h(\bfu, \bfm, \bfm) = 0$,
which is essential for the energy estimate. The extra $\half(\dv\bfu)$
term and the face average/jump splitting achieve this.
Let us verify: set $\bfz = \bfm$ in \eqref{eq:DG-trilinear}. The cell
terms give
$\sum_K \int_K [(\bfu\cdot\grad)\bfm]\cdot\bfm + \half(\dv\bfu)|\bfm|^2$
$= \sum_K \half\int_K \bfu\cdot\grad|\bfm|^2 + \half(\dv\bfu)|\bfm|^2$
$= \sum_K \half\int_K \dv(\bfu|\bfm|^2)$
$= \half\sum_K \int_{\partial K} (\bfu\cdot\bfn)|\bfm|^2$.
The face terms become
$-\half\sum_e \int_e (\bfu\cdot\bfn)\jump{\bfm}\cdot\avg{\bfm}
 -\half\sum_e \int_e (\bfu\cdot\bfn)\avg{\bfm}\cdot\jump{\bfm}$
$= -\sum_e \int_e (\bfu\cdot\bfn)\jump{\bfm}\cdot\avg{\bfm}$.
Using the identity $\jump{|\bfm|^2} = 2\avg{\bfm}\cdot\jump{\bfm}$
and $\sum_K \int_{\partial K}(\bfu\cdot\bfn)|\bfm|^2
= \sum_e \int_e (\bfu\cdot\bfn)\jump{|\bfm|^2}$,
everything cancels: $B_h(\bfu,\bfm,\bfm) = 0$. $\checkmark$
\end{derivnote}

\subsection{Upwind Stabilization}
\label{sec:upwind}

The skew-symmetric form $B_h$ is energy-neutral but does not suppress
oscillations. We add upwind stabilization:
\begin{equation}
  \label{eq:upwind}
  S_h(\bfu, \bfm, \bfz) :=
  \half \sum_{e \in \Eh} \int_e |\bfu \cdot \bfn|\,
  \jump{\bfm} \cdot \jump{\bfz} \dd s.
\end{equation}

The total transport form is:
\begin{equation}
  \label{eq:transport-total}
  \mathcal{B}_h(\bfu, \bfm, \bfz) = B_h(\bfu, \bfm, \bfz)
  + S_h(\bfu, \bfm, \bfz).
\end{equation}

\begin{derivnote}
Upwind adds dissipation: $S_h(\bfu,\bfm,\bfm) = \half\sum_e\int_e |\bfu\cdot\bfn|\,|\jump{\bfm}|^2 \geq 0$.
This is the DG equivalent of artificial viscosity, localized at element
interfaces where the solution has jumps. It vanishes for smooth solutions
(where $\jump{\bfm} \to 0$).
\end{derivnote}


% ══════════════════════════════════════════════════════════════════════
% SECTION 9: THE DISCRETE SCHEME
% ══════════════════════════════════════════════════════════════════════
\section{Energy-Stable Discrete Scheme}
\label{sec:scheme}

\subsection{Time Discretization}

Let $\delta t > 0$ be the time step, $t^n = n\,\delta t$. We use
backward Euler for time derivatives and treat the nonlinear terms
carefully to maintain energy stability.

Notation: $d_t v^n := \frac{v^n - v^{n-1}}{\delta t}$.

\subsection{The Coupled Scheme}

Given $(\theta^{n-1}, \bfu^{n-1}, \bfm^{n-1}, \varphi^{n-1})$, find
$(\theta^n, \psi^n, \bfu^n, p^n, \bfm^n, \varphi^n)$ such that:

\bigskip
\noindent\textbf{Block 1: Cahn--Hilliard.}
For all $(s, r) \in S_h \times S_h$:
\begin{subequations}
\label{eq:disc-CH}
\begin{align}
  (d_t\theta^n, s)
  + (\bfu^{n-1}\cdot\grad\theta^n, s)
  + \gamma(\grad\psi^n, \grad s)
  &= 0,
  \label{eq:disc-CH-a} \\
  (\psi^n, r) - \eps(\grad\theta^n, \grad r)
  - \frac{1}{\eps}\left(f(\theta^{n-1}) + f'(\theta^{n-1})(\theta^n - \theta^{n-1}), r\right)
  &= 0.
  \label{eq:disc-CH-b}
\end{align}
\end{subequations}

\begin{derivnote}
The nonlinearity $f(\theta)$ is linearized by a first-order Taylor
expansion about $\theta^{n-1}$. Since $|f'| \leq 2$
(from \eqref{eq:f-bounds}), this linearization is Lipschitz-stable.
The convex-concave splitting could also be used here.
\end{derivnote}

\bigskip
\noindent\textbf{Block 2: Navier--Stokes.}
For all $(\bfv, q) \in \bfV{}_h \times Q_h$:
\begin{subequations}
\label{eq:disc-NS}
\begin{align}
  (d_t\bfu^n, \bfv)
  + c_h(\bfu^{n-1}; \bfu^n, \bfv)
  + 2(\nu(\theta^n)\,\bm{D}(\bfu^n), \bm{D}(\bfv))
  &\notag\\
  - (p^n, \dv\bfv)
  &= \frac{\lambda}{\eps}(\theta^n\,\grad\psi^n, \bfv)
     + \mu_0\left((\bfm^{n-1}\cdot\grad)\bfh^{n-1}, \bfv\right)
     + (\bff_g^n, \bfv),
  \label{eq:disc-NS-a} \\
  (\dv\bfu^n, q) &= 0.
  \label{eq:disc-NS-b}
\end{align}
\end{subequations}
Here $c_h$ is the skew-symmetric CG convection form:
\begin{equation}
  \label{eq:cg-convection}
  c_h(\bm{w}; \bfu, \bfv) = ((\bm{w}\cdot\grad)\bfu, \bfv)
  + \half((\dv\bm{w})\bfu, \bfv).
\end{equation}

\begin{remark}
  The Kelvin force uses $\bfm^{n-1}$ and $\bfh^{n-1}$ (explicit in
  magnetization). This decouples NS from the magnetization-Poisson
  block but means the energy estimate requires care.
\end{remark}

\bigskip
\noindent\textbf{Block 3: Magnetization + Poisson (coupled, Picard).}
For all $(\bfz, \xi) \in \bfM_h \times X_h$:
\begin{subequations}
\label{eq:disc-mag}
\begin{align}
  (d_t\bfm^n, \bfz)
  + \mathcal{B}_h(\bfu^n, \bfm^n, \bfz)
  + \frac{1}{\tau}(\bfm^n, \bfz)
  &= \frac{\chi_0(\theta^n)}{\tau}(\bfh^n, \bfz),
  \label{eq:disc-mag-a} \\
  (\grad\varphi^n, \grad\xi) &= (\bfm^n, \grad\xi),
  \label{eq:disc-mag-b} \\
  \bfh^n &= \bfh_a - \grad\varphi^n.
  \label{eq:disc-mag-c}
\end{align}
\end{subequations}

\begin{keypoint}
Equations \eqref{eq:disc-mag-a}--\eqref{eq:disc-mag-c} are
\emph{nonlinearly coupled}: $\bfm^n$ depends on $\bfh^n$ (through the
RHS), and $\bfh^n$ depends on $\bfm^n$ (through the Poisson equation).
This is solved by \emph{Picard iteration} (Section~\ref{sec:picard}).
\end{keypoint}


% ══════════════════════════════════════════════════════════════════════
% SECTION 10: PICARD ITERATION
% ══════════════════════════════════════════════════════════════════════
\section{Picard Iteration for Magnetization--Poisson}
\label{sec:picard}

\subsection{The Coupling Problem}

The magnetization equation~\eqref{eq:disc-mag-a} has $\bfh^n$ on the
RHS, and $\bfh^n = \bfh_a - \grad\varphi^n$ where $\varphi^n$ depends
on $\bfm^n$. Writing out:
\[
  \frac{1}{\tau}(\bfm^n, \bfz)
  + \text{(other LHS terms)}
  = \frac{\chi_0}{\tau}(\bfh_a - \grad\varphi^n[\bfm^n], \bfz).
\]
This is nonlinear because $\varphi^n$ depends on $\bfm^n$ through the
Poisson equation.

\subsection{Block Gauss--Seidel (Picard) Algorithm}

Initialize: $\bfm^{n,0} = \bfm^{n-1}$, $\varphi^{n,0} = \varphi^{n-1}$.

For $k = 0, 1, 2, \ldots$ until convergence:
\begin{enumerate}
  \item \textbf{Magnetization sub-step:} Given $\varphi^{n,k}$, solve for
        $\bfm^{n,k+1}$:
        \begin{equation}
          \label{eq:picard-mag}
          (d_t\bfm^{n,k+1}, \bfz)
          + \mathcal{B}_h(\bfu^n, \bfm^{n,k+1}, \bfz)
          + \frac{1}{\tau}(\bfm^{n,k+1}, \bfz)
          = \frac{\chi_0}{\tau}(\bfh_a - \grad\varphi^{n,k}, \bfz).
        \end{equation}

  \item \textbf{Poisson sub-step:} Given $\bfm^{n,k+1}$, solve for
        $\varphi^{n,k+1}$:
        \begin{equation}
          \label{eq:picard-poisson}
          (\grad\varphi^{n,k+1}, \grad\xi) = (\bfm^{n,k+1}, \grad\xi)
          \qquad \forall\,\xi \in X_h.
        \end{equation}
\end{enumerate}

\begin{derivnote}
\textbf{Convergence:} The iteration is a contraction if
$\chi_0 < \chi_*$ for some critical value. Intuitively: the Poisson
equation maps $\bfm \mapsto \varphi$ with a ``gain'' related to the
Poincar\'e constant, and the magnetization equation maps
$\varphi \mapsto \bfm$ with gain $\chi_0/\tau$ divided by the
$(1/\tau + 1/\delta t)$ from the LHS. The product of gains must be $< 1$.

The relaxation parameter $\omega = 1/(1+\chi_0)$ can be used to
accelerate convergence by under-relaxing the Poisson update:
$\varphi^{n,k+1} \leftarrow \omega\,\varphi^{n,k+1} + (1-\omega)\varphi^{n,k}$.

In practice, convergence is fast: 3--7 iterations suffice with tolerance
$\sim 10^{-6}$.
\end{derivnote}


% ══════════════════════════════════════════════════════════════════════
% SECTION 11: SIMPLIFIED MODEL h = h_a
% ══════════════════════════════════════════════════════════════════════
\section{Simplified Model: \texorpdfstring{$\bfh = \bfh_a$}{h = h\_a} (No Demagnetization)}
\label{sec:simplified}

\subsection{When Is This Valid?}

If the demagnetizing field $\bfh_d = -\grad\varphi$ is negligible
compared to the applied field $\bfh_a$, we can set $\bfh = \bfh_a$ and
drop the Poisson equation entirely.

This is valid when:
\begin{itemize}[nosep]
  \item The ferrofluid region is thin (demagnetization is weak),
  \item $\chi_0$ is small (weak magnetization),
  \item The applied field is strong and slowly varying.
\end{itemize}

\subsection{Simplified System}

The magnetization equation becomes:
\begin{equation}
  \label{eq:simple-mag}
  \dt\bfm + (\bfu\cdot\grad)\bfm + \frac{1}{\tau}\bfm
  = \frac{\chi_0(\theta)}{\tau}\bfh_a.
\end{equation}
This is a \emph{linear} equation for $\bfm$ (no Picard needed).

\begin{remark}
  Convergence of the FEM is provable for this simplified model: optimal
  error estimates in $L^2$ and $H^1$ norms hold under standard regularity
  assumptions. The full model with Picard has convergence proven only
  under smallness assumptions on $\chi_0$.
\end{remark}


% ══════════════════════════════════════════════════════════════════════
% SECTION 12: ROSENSWEIG INSTABILITY
% ══════════════════════════════════════════════════════════════════════
\section{Rosensweig Instability: Benchmark Setup}
\label{sec:rosensweig}

\subsection{Physical Description}

The Rosensweig (normal-field) instability occurs when a uniform
vertical magnetic field is applied to a horizontal layer of ferrofluid.
Above a critical field strength, the flat surface becomes unstable and
forms a regular pattern of peaks (``spikes'').

The instability is a competition between:
\begin{itemize}
  \item \textbf{Magnetic force} (destabilizing): the Kelvin force
        $\mu_0(\bfm\cdot\grad)\bfh$ pulls ferrofluid upward where the
        field concentrates at a perturbation.
  \item \textbf{Gravity} (stabilizing): pulls the spike back down.
  \item \textbf{Surface tension} (stabilizing): penalizes curvature.
\end{itemize}

\subsection{Domain and Initial Conditions}

\paragraph{Uniform field case:}
\begin{itemize}[nosep]
  \item Domain: $\Omega = (0, L_x) \times (0, L_y)$ (2D).
  \item Ferrofluid layer on bottom, non-magnetic on top:
    \[
      \theta_0(x, y) = -\tanh\!\left(\frac{y - y_0}{\sqrt{2}\,\eps}\right),
    \]
    where $y_0$ is the unperturbed interface height.
  \item Applied field: $\bfh_a = H_a\,\bm{e}_y$ (uniform, vertical).
  \item $\bfu_0 = \bm{0}$, $\bfm_0 = \chi_0(\theta_0)\,\bfh_a$
    (equilibrium magnetization).
\end{itemize}

\subsection{Nondimensional Parameters}

The key nondimensional parameters are:
\begin{itemize}
  \item \textbf{Magnetic Bond number:}
    $\mathrm{Bo}_m = \frac{\mu_0 \chi_0^2 H_a^2 L}{\sigma}$
    (magnetic force / surface tension).
  \item \textbf{Gravitational Bond number:}
    $\mathrm{Bo}_g = \frac{\Delta\rho\,g\,L^2}{\sigma}$
    (gravity / surface tension).
  \item \textbf{Reynolds number:}
    $\mathrm{Re} = \frac{\rho\,U\,L}{\nu}$.
  \item \textbf{Cahn number:}
    $\mathrm{Cn} = \eps / L$ (interface width / domain size).
\end{itemize}

\subsection{What Should Happen}

\begin{enumerate}
  \item Initially, the flat interface is in equilibrium.
  \item Numerical perturbations (or explicit perturbation) seed the instability.
  \item If $\mathrm{Bo}_m$ exceeds a critical value, spikes grow.
  \item The spike pattern depends on the domain size and boundary conditions
        (periodic $\to$ regular array; bounded $\to$ fewer spikes near walls).
  \item In steady state, a fixed spike height is reached where magnetic
        force = gravity + surface tension.
\end{enumerate}

\begin{keypoint}
For the code to produce spikes, the following must all be correct:
\begin{enumerate}[nosep]
  \item The Kelvin force $\mu_0(\bfm\cdot\grad)\bfh$ must create an
        upward force where the interface bulges.
  \item The demagnetizing field (from Poisson) must concentrate field
        lines at the spike tips.
  \item The phase field must track the moving interface.
  \item Gravity must oppose the magnetic force.
  \item Surface tension must regularize the spike shape.
\end{enumerate}
If any one of these is wrong (sign error, missing term, wrong coupling),
the interface remains flat.
\end{keypoint}


% ══════════════════════════════════════════════════════════════════════
% SECTION 13: DIAGNOSTIC CHECKLIST
% ══════════════════════════════════════════════════════════════════════
\section{Diagnostic Checklist for Flat Interface Bug}
\label{sec:diagnostics}

Based on the derivations above, if the code produces a flat interface
when it should produce spikes, the possible causes are:

\subsection{Kelvin Force}
\begin{enumerate}[label=(\alph*)]
  \item Is $\mu_0(\bfm\cdot\grad)\bfh$ assembled correctly? Check sign
        and coefficient.
  \item Is the full field $\bfh = \bfh_a - \grad\varphi$ used, or only
        $\bfh_a$? Without the demagnetizing field, there is no field
        concentration at perturbations $\Rightarrow$ no instability.
  \item Is $(\bfm\cdot\grad)\bfh$ computed as
        $m_i\,\partial_i h_j$ or as $\grad(\bfm\cdot\bfh)$?
        They differ unless $\curl\bfh = 0$ \emph{and} $\curl\bfm = 0$.
\end{enumerate}

\subsection{Poisson / Demagnetization}
\begin{enumerate}[label=(\alph*)]
  \item Is the sign correct? $-\Delta\varphi = -\dv\bfm$ (note both
        negative signs).
  \item Is the Neumann BC $\dn\varphi = \bfm\cdot\bfn$ applied?
  \item Is the Picard iteration converging? Check residual.
  \item After Poisson solve, is $\bfh^n = \bfh_a - \grad\varphi^n$
        correctly computed?
\end{enumerate}

\subsection{Magnetization}
\begin{enumerate}[label=(\alph*)]
  \item Is the DG transport form correct? Check face integrals, signs,
        upwind direction.
  \item Is $\chi_0(\theta)$ correct? It should be $\approx \chi_{0,1}$
        in ferrofluid and $\approx 0$ in non-magnetic fluid.
  \item Is $\tau$ physical? Too small $\Rightarrow$ stiff, too large
        $\Rightarrow$ slow relaxation.
\end{enumerate}

\subsection{Cahn--Hilliard}
\begin{enumerate}[label=(\alph*)]
  \item Is the interface width $\eps$ resolved by the mesh? Need
        $h \lesssim \eps$.
  \item Is the mobility $\gamma$ reasonable? Too large $\Rightarrow$
        fast diffusion smears interface.
\end{enumerate}

\subsection{Energy}
\begin{enumerate}[label=(\alph*)]
  \item Is the total energy \eqref{eq:total-energy} decreasing (for
        zero applied field) or bounded (with applied field)?
  \item If energy grows without bound, there is an instability in the
        scheme — likely a sign error or missing stabilization.
\end{enumerate}


% ══════════════════════════════════════════════════════════════════════
\section*{Summary}

This document derives the complete two-phase ferrohydrodynamic system
from first principles:
\begin{itemize}[nosep]
  \item Cahn--Hilliard from Ginzburg--Landau free energy minimization,
  \item Simplified Shliomis from the full ferrohydrodynamic model,
  \item Poisson equation from Maxwell's magnetostatics,
  \item Kelvin force from magnetic dipole energy,
  \item Capillary force from the capillary stress tensor,
  \item Energy-stable FEM with DG magnetization and CG Poisson,
  \item DG transport form with upwind from element-wise integration by parts,
  \item Picard iteration from the $\bfm$--$\varphi$ coupling structure.
\end{itemize}

Every equation and every sign has a physical or mathematical justification.
The diagnostic checklist in Section~\ref{sec:diagnostics} provides a
systematic way to verify the code against this derivation.

\end{document}
